\chapter{The Binomial Model}\label{ch:dv:the-binomial-model}

The whiteboard shows a share at 100 and three steps. At each step the share
goes up 10\% or down 10\%, and a bank deposit grows by 2\%. The question is the
price of a call struck at 100 that expires after the third step. There is no
volatility number on the board, no probability of an up move and no formula,
and the answer, 10.3603, needs none of them. Every price in this chapter is the
cost of a portfolio of shares and deposit that pays exactly what the option
pays, node by node; the probability that appears on the way is the one
chapter 1 called risk-neutral, and it is an output of the prices, not an input.
The tree is the smallest \omterm{def:dv:no-arbitrage-and-the-fundamental-theorems:complete}{complete market} in which options are non-trivial,
and it is also a working pricer: with enough steps it converges to the
Black--Scholes value of chapter 3, and it prices American exercise, which no
closed formula does.

\section{One period: replication and the risk-neutral probability}

\begin{definition}[Binomial model]\label{def:dv:the-binomial-model:binomial}
In the \emph{binomial model}\index{binomial model} time runs in $n$ steps of
length $\Delta t=T/n$. Over each step the share price is multiplied by $u$ or by
$d<u$, and one unit of cash deposited grows to $R=e^{r\Delta t}$ (with a
continuous dividend yield $q$, the share's own growth is compared with
$e^{(r-q)\Delta t}$). The model has no arbitrage if and only if
$d<e^{(r-q)\Delta t}<u$.
\end{definition}

Over one step, a claim paying $V_u$ after an up move and $V_d$ after a down move
is replicated by $\Delta$ shares and a deposit $b$ if $\Delta S u+bR=V_u$ and
$\Delta Sd+bR=V_d$. Two equations, two unknowns:
\begin{equation}\label{eq:dv:the-binomial-model:onestep}
\Delta=\frac{V_u-V_d}{S(u-d)},\qquad
V=\Delta S+b=\frac{1}{R}\bigl(pV_u+(1-p)V_d\bigr),\qquad p=\frac{R-d}{u-d}.
\end{equation}
The weight $p$ lies in $(0,1)$ exactly under the no-arbitrage condition; it is
the risk-neutral probability of the up move, the normalised \omterm{def:dv:no-arbitrage-and-the-fundamental-theorems:ad}{state price} of
chapter 1, and the real-world probability of an up move appears nowhere. Two
traders who disagree about that probability agree on the price, because each can
build the option from shares and cash.

\begin{example}[One step]\label{ex:dv:the-binomial-model:onestep}
A share at 50 will be at 60 or 40; rates are zero. A call struck at 50 pays 10 or
0: $\Delta=10/20=0.5$ share, and the deposit is $b=(0-0.5\times40)=-20$, a
borrowing of 20. The call costs $0.5\times50-20=5$, and $p=(1-0.8)/0.4=0.5$.
\end{example}

\section{Many periods: the recombining tree}

\begin{definition}[Recombining tree]\label{def:dv:the-binomial-model:recombining}
A \emph{recombining tree}\index{recombining tree} is a tree in which an up move
followed by a down move reaches the same node as the reverse: after $k$ steps
the share is at one of the $k+1$ values $Su^jd^{k-j}$, $j=0,\dots,k$, instead of
$2^k$. Pricing on it costs $O(n^2)$ operations and $O(n)$ memory.
\end{definition}

Applying \eqref{eq:dv:the-binomial-model:onestep} node by node from the last
step back to the first is the backward induction of dynamic programming (One
Quant Book 4, chapter 9). The portfolio at each node is rebalanced at the next
using only its own value: it is a \omterm{def:dv:no-arbitrage-and-the-fundamental-theorems:sf}{self-financing strategy} in the sense of
chapter 1, and it replicates the claim on every path.

\begin{proposition}[The binomial price]\label{prop:dv:the-binomial-model:formula}
A European claim paying $g(S_T)$ after $n$ steps is worth
\[
V_0=R^{-n}\sum_{j=0}^{n}\binom{n}{j}p^j(1-p)^{n-j}\,g\bigl(S_0u^jd^{n-j}\bigr),
\]
the discounted expectation of the payoff when the number of up moves is binomial
with parameters $n$ and $p$.
\end{proposition}
\begin{proof}
By induction on the number of remaining steps: each backward step replaces two
values by their $p$-weighted average divided by $R$, and the number of paths
reaching node $j$ is $\binom{n}{j}$.
\end{proof}

\begin{omfigure}
\begin{tikzpicture}[font=\footnotesize,x=2.9cm,y=0.95cm,
  nd/.style={draw=omDef,thick,rounded corners=2pt,fill=omDef!6,align=center,inner sep=2pt,minimum width=1.55cm},
  ex/.style={nd,draw=omThm,fill=omThm!8}]
\node[nd] (a) at (0,0) {100\\4.908};
\node[ex] (b0) at (1,-1.5) {90\\\textbf{10.000}};
\node[nd] (b1) at (1,1.5) {110\\1.676};
\node[ex] (c0) at (2,-3) {81\\\textbf{19.000}};
\node[nd] (c1) at (2,0) {99\\4.275};
\node[nd] (c2) at (2,3) {121\\0.000};
\node[nd] (d0) at (3,-4.5) {72.9\\27.100};
\node[nd] (d1) at (3,-1.5) {89.1\\10.900};
\node[nd] (d2) at (3,1.5) {108.9\\0.000};
\node[nd] (d3) at (3,4.5) {133.1\\0.000};
\foreach \s/\t in {a/b0,a/b1,b0/c0,b0/c1,b1/c1,b1/c2,c0/d0,c0/d1,c1/d1,c1/d2,c2/d2,c2/d3}
  \draw[gray,-{Stealth}] (\s) -- (\t);
\node[anchor=west,align=left,text=omThm] at (0.1,-3.9) {red: exercise now\\beats holding};
\end{tikzpicture}

\omcaption{The whiteboard tree: share price (top line) and American put value
(bottom line) at every node, $u=1.1$, $d=0.9$, 2\% per step, strike 100. At 90
and at 81 the put is worth more exercised (10 and 19) than held (9.196 and
17.039), and the tree takes the larger value. Data: the tutorial.}\label{fig:dv:the-binomial-model:tree}
\end{omfigure}

\begin{example}[The whiteboard call]\label{ex:dv:the-binomial-model:whiteboard}
With $u=1.1$, $d=0.9$, $R=1.02$: $p=(1.02-0.9)/0.2=0.6$. After three steps the
share is at 133.1, 108.9, 89.1 or 72.9, and the call pays 33.1, 8.9, 0 or 0. Its
value is $(0.6^3\times33.1+3\times0.6^2\times0.4\times8.9)/1.02^3=10.3603$. At
the first node the \omterm{def:dv:no-arbitrage-and-the-fundamental-theorems:claim}{replicating portfolio} holds $\Delta=0.6240$ share and a
deposit of $-52.04$; put--call parity holds on the tree, $C-P=100-100/1.02^3=5.7678$.
\end{example}

\section{Calibrating the tree to a volatility}

To model a share with volatility $\sigma$ the step factors must reproduce, to first
order in $\Delta t$, the mean and variance of the log-return over a step. Two
choices are standard; a third is built for accuracy.

\begin{definition}[Cox--Ross--Rubinstein and Jarrow--Rudd trees]\label{def:dv:the-binomial-model:crr}
The \emph{Cox--Ross--Rubinstein tree}\index{Cox--Ross--Rubinstein tree} takes
$u=e^{\sigma\sqrt{\Delta t}}$ and $d=1/u$, so that the tree is symmetric in
log-price and the drift enters only through $p$. The \emph{Jarrow--Rudd
tree}\index{Jarrow--Rudd tree} takes
$u,d=\exp\bigl((r-q-\tfrac12\sigma^2)\Delta t\pm\sigma\sqrt{\Delta t}\bigr)$, which
puts the drift into the factors and makes $p$ close to one half.
\end{definition}

\begin{proposition}[Moment matching]\label{prop:dv:the-binomial-model:moments}
In both trees the one-step log-return has risk-neutral mean
$(r-q-\tfrac12\sigma^2)\Delta t+O(\Delta t^{3/2})$ and variance
$\sigma^2\Delta t+O(\Delta t^2)$, and the growth of the share matches
$e^{(r-q)\Delta t}$ exactly by the choice of $p$.
\end{proposition}
\begin{proof}
For the \omterm{def:dv:the-binomial-model:crr}{Cox--Ross--Rubinstein tree}, expand
$p=(e^{(r-q)\Delta t}-e^{-\sigma\sqrt{\Delta t}})/(2\sinh\sigma\sqrt{\Delta t})=
\tfrac12+\tfrac{(r-q-\sigma^2/2)\sqrt{\Delta t}}{2\sigma}+O(\Delta t)$; the mean of
$\pm\sigma\sqrt{\Delta t}$ is $(2p-1)\sigma\sqrt{\Delta t}$ and its second moment
$\sigma^2\Delta t$. The Jarrow--Rudd computation is the same with the drift
moved into the factors.
\end{proof}

By the central limit theorem (One Quant Book 4, chapter 1) the number of up moves,
centred and scaled, becomes normal: the terminal log-price of the tree converges
in law to the normal distribution of geometric Brownian motion
(\cref{fig:dv:the-binomial-model:terminal}), and the tree price of a European
option converges to the Black--Scholes price of chapter 3.

\begin{omfigure}
\begin{tikzpicture}
\begin{axis}[width=10.5cm,height=4.8cm,ybar,bar width=4pt,
  xlabel={log-return at expiry},ylabel={density},
  tick label style={font=\small},label style={font=\small},
  xmin=-0.85,xmax=0.85,ymin=0,ymax=2.4,ymajorgrids,grid style={gray!25},
  legend columns=2,legend style={font=\footnotesize,at={(0.5,-0.3)},anchor=north,draw=none,fill=none,/tikz/every even column/.append style={column sep=8pt}},
  table/col sep=comma]
\addplot[fill=omDef!55,draw=omDef] table[x=x,y=tree]{figdata/derivatives/02-the-binomial-model/terminal.csv};
\addplot[omThm,very thick,sharp plot,mark=none] table[x=x,y=normal]{figdata/derivatives/02-the-binomial-model/terminal.csv};
\legend{50-step tree (probability / spacing),normal limit}
\end{axis}
\end{tikzpicture}

\omcaption{The risk-neutral law of the terminal log-return in a 50-step
\omterm{def:dv:the-binomial-model:crr}{Cox--Ross--Rubinstein tree} (one year, $r=5\%$, $\sigma=20\%$), each probability
divided by the spacing $2\sigma\sqrt{\Delta t}$ between nodes, against the normal
density with mean $(r-\sigma^2/2)T$ and standard deviation $\sigma\sqrt T$. Data:
the tutorial.}\label{fig:dv:the-binomial-model:terminal}
\end{omfigure}

\section{Convergence and its oscillations}

For a one-year at-the-money put with $r=5\%$ and $\sigma=20\%$ the Black--Scholes
value is 5.5735. The Cox--Ross--Rubinstein price with 100 steps is 0.0200 below
it, with 101 steps 0.0174 above: the error falls like $1/n$ and changes sign
with the parity of $n$ (\cref{fig:dv:the-binomial-model:conv}). The cause is the
payoff's kink. The nodes near the strike move relative to it as $n$ changes,
and the piecewise-linear payoff is sampled at a different point of its kink
each time; with the strike exactly at the spot, even $n$ puts a node on the kink
and odd $n$ straddles it.

\begin{definition}[Leisen--Reimer tree]\label{def:dv:the-binomial-model:lr}
The \emph{Leisen--Reimer tree}\index{Leisen--Reimer tree} uses an odd number of
steps and chooses $p$ and $u$ so that the tree's probabilities of finishing above
the strike, under the risk-neutral and the share measures, equal $\Phi(d_2)$ and
$\Phi(d_1)$ (through a normal-to-binomial inversion due to Peizer and Pratt).
The tree is centred on the strike, and the error for European options falls
like $1/n^2$ without oscillation.
\end{definition}

\begin{omfigure}
\begin{tikzpicture}
\begin{axis}[width=11cm,height=5cm,
  xlabel={number of steps $n$},ylabel={tree price $-$ formula},
  tick label style={font=\small},label style={font=\small},
  xmin=10,xmax=200,ymin=-0.1,ymax=0.1,grid=major,grid style={gray!25},
  yticklabel style={/pgf/number format/.cd,fixed,precision=2},
  legend columns=2,legend style={font=\footnotesize,at={(0.5,-0.3)},anchor=north,draw=none,fill=none,/tikz/every even column/.append style={column sep=8pt}},
  table/col sep=comma]
\addplot[omDef,thick] table[x=n,y=crr]{figdata/derivatives/02-the-binomial-model/conv_even.csv};
\addplot[omDef,thick,dashed] table[x=n,y=crr]{figdata/derivatives/02-the-binomial-model/conv_odd.csv};
\addplot[omProp,thick] table[x=n,y=jr]{figdata/derivatives/02-the-binomial-model/conv_even.csv};
\addplot[omThm,very thick] table[x=n,y=lr]{figdata/derivatives/02-the-binomial-model/conv_lr.csv};
\legend{Cox--Ross--Rubinstein even $n$,Cox--Ross--Rubinstein odd $n$,Jarrow--Rudd even $n$,Leisen--Reimer odd $n$}
\end{axis}
\end{tikzpicture}

\omcaption{Error of three trees for the one-year at-the-money European put
($r=5\%$, $\sigma=20\%$, formula 5.5735), even and odd step counts drawn
separately. The Cox--Ross--Rubinstein error is negative for even $n$, positive
for odd $n$, and shrinks like $1/n$; the
Leisen--Reimer error is invisible at this scale from 11 steps. Data: the
tutorial.}\label{fig:dv:the-binomial-model:conv}
\end{omfigure}

\begin{remark}[What the oscillation costs]\label{rem:dv:the-binomial-model:cost}
The Cox--Ross--Rubinstein put stays within one cent of the formula only from 199
steps on; the \omterm{def:dv:the-binomial-model:lr}{Leisen--Reimer tree} is within 0.3 cent at 11 steps and within
$10^{-4}$ at 101. Averaging the prices at $n$ and $n+1$ cancels most of the
oscillation at the cost of a second tree; Richardson extrapolation (One Quant
Book 4, chapter 27) removes the $1/n$ term once the oscillation is gone. Chapter
22 compares trees with finite-difference grids on accuracy against run time.
\end{remark}

\section{American exercise on a tree}

An American option can be exercised at any node. At each node the holder compares
the value of continuing, given by \eqref{eq:dv:the-binomial-model:onestep}, with
the value of exercising now, and the option is worth the larger:
\[
V_k^j=\max\Bigl(g(S_k^j),\ R^{-1}\bigl(pV_{k+1}^{j+1}+(1-p)V_{k+1}^{j}\bigr)\Bigr).
\]
The writer's \omterm{def:dv:no-arbitrage-and-the-fundamental-theorems:claim}{replicating portfolio} is still self-financing where the holder
continues; where the holder exercises, it is liquidated to pay the exercise
value. If a holder fails to exercise at a node where exercise is optimal, the
writer's portfolio is worth more than the option from then on: the writer
profits from the holder's mistake, a theme of chapter 6.

\begin{proposition}[No early exercise of a call without dividends]\label{prop:dv:the-binomial-model:call}
On a tree with $r\ge0$ and $q=0$ an American call is worth the European call:
exercise is never strictly optimal before expiry.
\end{proposition}
\begin{proof}
At any node the continuation value is at least
$R^{-1}\bigl(pS_{k+1}^{j+1}+(1-p)S_{k+1}^{j}\bigr)-KR^{-(n-k)}$ by convexity of the
payoff and induction, which equals $S_k^j-KR^{-(n-k)}\ge S_k^j-K$: the call is
worth more alive than exercised.
\end{proof}

A put has no such bound: deep in the money its holder would rather have the
strike now and earn interest on it (\cref{fig:dv:the-binomial-model:american}).
On the whiteboard tree the American put is worth 4.9076 against 4.5925 for the
European: an early-exercise premium of 0.3151, earned at the two nodes of
\cref{fig:dv:the-binomial-model:tree} where exercise wins.

\begin{omfigure}
\begin{tikzpicture}
\begin{axis}[width=10.5cm,height=5cm,
  xlabel={spot},ylabel={put value},
  tick label style={font=\small},label style={font=\small},
  xmin=60,xmax=140,ymin=0,ymax=42,grid=major,grid style={gray!25},
  legend columns=3,legend style={font=\footnotesize,at={(0.5,-0.3)},anchor=north,draw=none,fill=none,/tikz/every even column/.append style={column sep=8pt}},
  table/col sep=comma]
\addplot[gray,thick,dashed] table[x=spot,y=intrinsic]{figdata/derivatives/02-the-binomial-model/american.csv};
\addplot[omDef,very thick] table[x=spot,y=european]{figdata/derivatives/02-the-binomial-model/american.csv};
\addplot[omThm,very thick] table[x=spot,y=american]{figdata/derivatives/02-the-binomial-model/american.csv};
\legend{exercise value $(K-S)^+$,European,American}
\end{axis}
\end{tikzpicture}

\omcaption{One-year puts struck at 100 ($r=5\%$, $\sigma=20\%$, 501-step
\omterm{def:dv:the-binomial-model:lr}{Leisen--Reimer tree}). The European put falls below its exercise value deep in the
money, where waiting costs the interest on the strike; the American put never
does, and meets the exercise value tangentially at the exercise boundary, near
81.1 here. Data: the tutorial.}\label{fig:dv:the-binomial-model:american}
\end{omfigure}

\section{Tutorial: a tree pricer and its convergence}\label{tut:dv:the-binomial-model:tut}

\begin{tutorial}
\textbf{Goal.} Build the three trees, price Europeans and Americans on them, and
measure their convergence to the formula. \textbf{End state:}
\cref{fig:dv:the-binomial-model:tree,fig:dv:the-binomial-model:conv,fig:dv:the-binomial-model:american}
and the numbers of \cref{ex:dv:the-binomial-model:whiteboard}.
\begin{tutsteps}
\item \textbf{Step factors.} One function returns $(u,d,p)$ for each tree and
refuses a tree with an arbitrage.
\omcode{code/firm/binomial/firm_binomial.py}{21}{46}{Step factors and risk-neutral probability of the three trees.}\label{lst:dv:the-binomial-model:params}
\item \textbf{Backward induction.} One array per time slice, shortened by one
at each step; American exercise is one \texttt{np.maximum} per slice. The
dividend argument is used in chapter 5.
\omcode{code/firm/binomial/firm_binomial.py}{49}{70}{Vectorised backward induction, European or American.}\label{lst:dv:the-binomial-model:price}
\item \textbf{Run} \texttt{dv\_binomial.whiteboard()} for the node values of
\cref{fig:dv:the-binomial-model:tree}, then \texttt{error(n, method)} for
$n=10,\dots,200$ and \texttt{fig\_binomial.py}.
\end{tutsteps}
\textbf{What to change next.} Put the strike at 101 instead of 100 and watch the
oscillation's phase move; then price the American put with the average of the
$n$- and $(n+1)$-step \omterm{def:dv:the-binomial-model:crr}{Cox--Ross--Rubinstein trees} and compare with Leisen--Reimer.
\end{tutorial}

\section{Build: the tree pricer}\label{bld:dv:the-binomial-model:binomial}

\begin{build}
\bfield{Purpose} The miniature firm's first pricer of American options, and the
reference against which the finite-difference engine of chapter 22 is tested.
\bfield{Interface} \texttt{params(method, spot, strike, t, r, q, vol, n) ->
(u, d, p)} with method in crr, jr, lr; \texttt{price(spot, strike, t, r, vol, n,
right, american, q, method, dividends)}; \texttt{node\_values(\ldots)} for small
trees.
\bfield{Rules} Continuous compounding; $n$ odd for Leisen--Reimer; an arbitrage in
the factors raises an error rather than returning a number; $O(n)$ memory.
\bfield{Acceptance tests} \texttt{code/firm/binomial/tests/}: parity on the tree;
convergence of all three trees to the formula; Leisen--Reimer within $10^{-4}$ at
101 steps; American call equal to European without dividends; American put above
European and above its exercise value.
\bfield{Stretch} A trinomial tree (chapter 22) and a pricer that returns delta and
gamma from the first two time slices at no extra cost.
\end{build}

\begin{omsources}
\item J.~C.~Cox, S.~A.~Ross and M.~Rubinstein, ``Option pricing: a simplified
approach'', \emph{Journal of Financial Economics} 7 (1979) 229--263.
\item R.~J.~Rendleman and B.~J.~Bartter, ``Two-state option pricing'',
\emph{Journal of Finance} 34 (1979) 1093--1110.
\item R.~A.~Jarrow and A.~Rudd, \emph{Option Pricing}, Irwin, 1983.
\item D.~P.~J.~Leisen and M.~Reimer, ``Binomial models for option valuation:
examining and improving convergence'', \emph{Applied Mathematical Finance} 3
(1996) 319--346.
\end{omsources}

\section{Exercises}

\begin{exercise}[$\star$]\label{exo:dv:the-binomial-model:1}
A share at 50 will be at 60 or 40; rates are zero. Give the \omterm{def:dv:no-arbitrage-and-the-fundamental-theorems:claim}{replicating portfolio}
and the price of a call struck at 50, and the risk-neutral probability.
\end{exercise}

\begin{exercise}[$\star$]\label{exo:dv:the-binomial-model:2}
Give $u$, $d$ and $p$ of a one-month Cox--Ross--Rubinstein step with $\sigma=20\%$
and $r=3\%$.
\end{exercise}

\begin{exercise}[$\star$]\label{exo:dv:the-binomial-model:3}
Check put--call parity on the whiteboard tree.
\end{exercise}

\begin{exercise}[$\star\star$]\label{exo:dv:the-binomial-model:4}
A colleague estimates that the share of the whiteboard goes up with probability
0.8, not 0.6, and says the call is therefore worth more than 10.3603. Explain why
the price does not depend on that estimate.
\end{exercise}

\begin{exercise}[$\star\star$]\label{exo:dv:the-binomial-model:5}
Price on the whiteboard tree a digital that pays 1 if the share ends above 100.
\end{exercise}

\begin{exercise}[$\star\star$]\label{exo:dv:the-binomial-model:6}
Compare the Cox--Ross--Rubinstein and Jarrow--Rudd factors and probabilities for
$\sigma=20\%$, $r=5\%$, $\Delta t=0.1$.
\end{exercise}

\begin{exercise}[$\star\star\star$]\label{exo:dv:the-binomial-model:7}
\emph{Coding.} Price the one-year at-the-money European put ($r=5\%$,
$\sigma=20\%$) with 101 steps on the Cox--Ross--Rubinstein and \omterm{def:dv:the-binomial-model:lr}{Leisen--Reimer
trees} and compare both with 5.5735.
\end{exercise}

\begin{exercise}[$\star\star\star$]\label{exo:dv:the-binomial-model:8}
\emph{Find the flaw.} ``We calibrated the tree to the data: up 5\%, down 3\%, and
6\% growth of cash per step. The probability comes out at 1.125, but we price
with it anyway.''
\end{exercise}

\section{Problem: The Whiteboard Tree}

\begin{problem}[{Weekend problem --- from three steps to a working pricer}]\label{pb:dv:the-binomial-model:1}
A candidate is given the whiteboard tree of this chapter (share 100, up 10\%,
down 10\%, cash growth 2\% per step, three steps, strike 100), then asked to turn
it into a pricer for a one-year option with $r=5\%$ and $\sigma=20\%$.

\medskip\noindent\textbf{Part I --- The European options.}
\begin{enumerate}
\item Give the risk-neutral probability of an up move.
\item List the share prices after three steps.
\item Price the call.
\item Give the \omterm{def:dv:no-arbitrage-and-the-fundamental-theorems:claim}{replicating portfolio} at the first node.
\item Price the European put and check put--call parity.
\end{enumerate}

\medskip\noindent\textbf{Part II --- American exercise.}
\begin{enumerate}[resume]
\item At which nodes is early exercise of the put optimal?
\item Price the American put.
\item Give the early-exercise premium.
\item Is the American call worth more than the European call on this tree?
\item Why does the holder at the node 90 prefer exercising?
\end{enumerate}

\medskip\noindent\textbf{Part III --- A real tree.}
\begin{enumerate}[resume]
\item Give the Black--Scholes value of the one-year at-the-money put.
\item Give the Cox--Ross--Rubinstein errors at 100 and 101 steps.
\item Give the Leisen--Reimer error at 101 steps.
\item Price the American put with a 2\,001-step \omterm{def:dv:the-binomial-model:lr}{Leisen--Reimer tree}, and its premium.
\item Why does the Cox--Ross--Rubinstein error change sign with the parity of $n$?
\end{enumerate}

\medskip\noindent\textbf{Part IV --- Judgement.}
\begin{enumerate}[resume]
\item How many steps does each tree need for one-cent accuracy?
\item How does the run time grow with $n$, and what does that imply for a desk
repricing 10\,000 American options?
\item What does a tree assume that a real market does not offer?
\item State the \emph{named result}: the early-exercise premium on the whiteboard
tree, and the number of Cox--Ross--Rubinstein steps after which the put stays
within one cent of the formula.
\item In one sentence: why is the tree's probability not a forecast?
\end{enumerate}
\end{problem}

\section{Interview questions}

\begin{interviewq}[$\star$ \iqroles{trader, researcher}]\label{iq:dv:the-binomial-model:1}
A share at 100 goes to 120 or 80 in one period, rates are zero. Price a call
struck at 100 and give the hedge.
\end{interviewq}

\begin{interviewq}[$\star$ \iqroles{researcher}]\label{iq:dv:the-binomial-model:2}
Why does the real-world probability of an up move not enter the binomial price?
\end{interviewq}

\begin{interviewq}[$\star\star$ \iqroles{researcher, developer}]\label{iq:dv:the-binomial-model:3}
Why does the Cox--Ross--Rubinstein price oscillate as the number of steps grows,
and how do you fix it?
\end{interviewq}

\begin{interviewq}[$\star\star$ \iqroles{trader}]\label{iq:dv:the-binomial-model:4}
Why might you exercise an American put early, but never an American call on a
share that pays no dividend?
\end{interviewq}

\begin{interviewq}[$\star\star$ \iqroles{developer}]\label{iq:dv:the-binomial-model:5}
How many nodes does an $n$-step \omterm{def:dv:the-binomial-model:recombining}{recombining tree} have, and how much memory does
pricing on it need?
\end{interviewq}

\begin{interviewq}[$\star\star\star$ \iqroles{developer, researcher}]\label{iq:dv:the-binomial-model:6}
Your tree pricer must return delta and gamma as well as the price, without
repricing. How?
\end{interviewq}
